% ==================================================================
% id:       prob-011
% topic:    מבוא להסתברות
% subtopic: עקרון ההכלה וההפרדה
% solution: full
% ==================================================================
מדען מבצע שני ניסויים. הסיכוי להצליח בניסוי הראשון הוא $0.4$, הסיכוי להצליח בשני הוא $0.3$, והסיכוי להצליח \emph{רק} בניסוי הראשון הוא $0.3$.
\begin{enumerate}
\item מה הסיכוי שיצליח רק בניסוי השני?
\begin{solbox}
נסמן $A$ = „הצלחה בראשון”, $B$ = „הצלחה בשני”. נתון $\PP(A\cap B^C)=0.3$, ולכן $\PP(A\cap B)=\PP(A)-\PP(A\cap B^C)=0.4-0.3=0.1$. מכאן $\PP(A^C\cap B)=\PP(B)-\PP(A\cap B)=0.3-0.1=\ans{0.2}$.
\end{solbox}
\item מה הסיכוי שייכשל לפחות בניסוי אחד?
\begin{solbox}
זה המשלים של „הצליח בשניהם”: $1-0.1=\ans{0.9}$.
\end{solbox}
\item מה הסיכוי שייכשל בדיוק בניסוי אחד?
\begin{solbox}
להיכשל בדיוק באחד זה להצליח בדיוק באחד: $\PP(A\cap B^C)+\PP(A^C\cap B)=0.3+0.2=\ans{0.5}$.
\end{solbox}
\end{enumerate}
